% Shear beam — analytical derivation
% AIM Interactive Mechanics series. Compile with pdflatex.
\documentclass[10pt,a4paper]{article}
\usepackage[margin=2cm]{geometry}
\usepackage{amsmath}
\usepackage[T1]{fontenc}
\usepackage{lmodern}
\usepackage{microtype}
\usepackage{titlesec}
\titlespacing*{\section}{0pt}{9pt}{2pt}
\usepackage{xcolor}
\definecolor{iron}{HTML}{935636}
\definecolor{celdark}{HTML}{4B6C61}
\definecolor{celline}{HTML}{CDE0D8}
\usepackage[colorlinks=true,urlcolor=iron,linkcolor=black]{hyperref}
\setlength{\parindent}{0pt}
\setlength{\parskip}{4pt}
\AtBeginDocument{\setlength{\abovedisplayskip}{6pt}\setlength{\belowdisplayskip}{6pt}\setlength{\abovedisplayshortskip}{3pt}\setlength{\belowdisplayshortskip}{3pt}}
\pagestyle{empty}

\begin{document}

{\footnotesize\textcolor{celdark}{\textsc{aim interactive mechanics \,\textperiodcentered\, analytical derivation}}}

{\LARGE The Shear Beam}

\textcolor{celdark}{\small Companion to the interactive simulator at
\url{https://aimlab.uk/sims/shear-beam.html}}

\medskip
{\color{celline}\hrule height 1pt}

\section*{Setup and assumptions}
A cantilever of length $L$ and rectangular cross-section $w \times h$ carries a
transverse tip load $F$; the material has shear modulus $G$. The idealisation is
opposite to Euler--Bernoulli: \emph{no bending stiffness}, deformation by shear
alone --- cross-sections translate laterally by $u(x)$ without rotating.

\section*{Kinematics}
With sections sliding but not rotating, the change of angle between the beam axis
and the cross-section is carried entirely by shear:
$\gamma = \mathrm{d}u/\mathrm{d}x$, taken as uniform over the section in this
idealisation.

\section*{Constitutive law and the correction factor}
The shear force is the resultant of the shear stress over the section. In
reality $\tau$ is not uniform --- for a rectangle it is parabolic,
\begin{equation}
\tau(y) \;=\; \frac{3V}{2A}\left(1 - \Big(\frac{2y}{h}\Big)^{\!2}\right),
\qquad
\tau_{\max} \;=\; \frac{3V}{2A} \quad \text{at mid-depth},
\end{equation}
vanishing on the free surfaces. Matching the shear strain energy of this
distribution to that of the uniform model gives the correction factor $k = 5/6$
for a rectangle. The section law is then
\begin{equation}
V \;=\; k G A \,\gamma \;=\; k G A \,u',
\qquad
A = w h .
\end{equation}
$kGA$ is the \emph{shear rigidity} --- it plays the role that $EI$ plays in
bending.

\section*{Governing equation}
Transverse equilibrium of a slice gives $\mathrm{d}V/\mathrm{d}x = -q = 0$
between loads, hence a second-order equation in $G$ (compare the fourth-order
$EI\,u'''' = 0$ of the bending beam):
\begin{equation}
k G A \,\frac{\mathrm{d}^2 u}{\mathrm{d}x^2} \;=\; 0,
\qquad
u(0) = 0, \qquad k G A\, u'(L) = F .
\end{equation}

\section*{Solution}
\begin{equation}
\boxed{\; u(x) \;=\; \frac{F x}{k G A} \;}
\qquad\Rightarrow\qquad
\delta \;=\; \frac{F L}{k G A},
\qquad
\gamma \;=\; \frac{F}{k G A} \;\; \text{(constant)} .
\end{equation}
The deflected shape is a straight line with a kink at the wall.

\section*{When does shear matter?}
For the same beam, the Euler--Bernoulli tip deflection is $F L^3 / 3EI$. The
ratio of the two is
\begin{equation}
\frac{\delta_{\text{shear}}}{\delta_{\text{bend}}}
\;=\; \frac{3 E I}{k G A L^2}
\;=\; \frac{E}{4 k G}\left(\frac{h}{L}\right)^{\!2},
\end{equation}
which scales with $(h/L)^2$: negligible for slender beams, significant for
stubby ones. Physical beams do both, which is captured by the Timoshenko beam formulation.

\section*{Worked example}
$G = 1$~MPa, $L = 500$~mm, $w = 25$~mm, $h = 50$~mm, $F = 60$~N:
\begin{equation}
A = 1250~\text{mm}^2, \qquad
kGA = \tfrac{5}{6} \times 10^6 \times 1.25\times10^{-3} = 1042~\text{N}, \qquad
\delta = \frac{60 \times 0.5}{1042} = 28.8~\text{mm},
\end{equation}
with $\gamma = 5.8\%$ and $\tau_{\max} = 3F/2A = 72$~kPa. Set the simulator
sliders to these values and compare.

\vfill
{\color{celline}\hrule height 0.6pt}
\smallskip
{\footnotesize\textcolor{celdark}{Part of the multimodal mechanics series
(analytical \textperiodcentered\ simulated \textperiodcentered\ experimental)
\,\textperiodcentered\, AIM Group, Imperial College London.}}

\end{document}
